\begin{ouverture}[college1]{7}{La matière dans l’Univers}{Partie 1 · Organisation et transformations de la matière}{De quoi les étoiles et la Terre sont-elles faites ?}
\objectif{Situer Terre, Système solaire et galaxie.}
\objectif{Comparer des échelles de distance et de temps.}
\objectif{Relier les éléments de la Terre à l’histoire de la matière.}
\prerequis{Proportionnalité et lecture des symboles chimiques.}
\end{ouverture}
\pcPage{Le cours}{La matière dans l’Univers}
\phantomsection\label{ess-7}
\essSection{1}{Des objets emboîtés}
La Terre est une planète du Système solaire ; le Soleil est son étoile. Le Système solaire appartient à la Voie lactée, une galaxie contenant un très grand nombre d’étoiles. L’Univers observable contient de nombreuses galaxies. Une planète ne produit pas sa lumière visible comme une étoile : elle peut diffuser la lumière reçue. Un satellite naturel tourne autour d’une planète dans le modèle étudié.\par
\pcFigure{figures/c08-systeme.pdf}{Ordre des huit planètes : les tailles et les espacements ne sont pas à l’échelle.}
\essSection{2}{Comparer des distances}
Pour le voisinage terrestre, on utilise souvent le kilomètre. L’unité astronomique, notée ua, vaut ici environ $150$ millions de kilomètres et correspond à l’échelle de la distance Terre-Soleil. L’année-lumière est une distance parcourue par la lumière en un an, environ $9{,}46\times10^{12}\unite{km}$. Ce n’est pas une durée. La vitesse de propagation sera étudiée au chapitre 16.\par
\clearpage
\essSection{3}{Une histoire de la matière}
Le Système solaire s’est formé il y a environ $4{,}6$ milliards d’années à partir de gaz et de poussières. L’Univers a évolué avant sa formation : les premiers éléments légers comprennent principalement hydrogène et hélium. Des étoiles et des événements stellaires ont ensuite produit et dispersé de nombreux éléments plus lourds. Carbone, oxygène, fer ou silicium ne sont donc pas propres à la Terre. Le modèle cosmologique situe l’âge de l’Univers vers 13,8 milliards d’années ; cet âge ne se confond pas avec celui de la Terre. Une frise doit distinguer ces échelles de temps.\par
\pcFigure{figures/c08-echelles.pdf}{De la planète à la galaxie : comparer des ordres de grandeur exige de préciser l’unité et le type de distance.}
\essSection{4}{Des ressources inégalement réparties}
La matière observable étudiée dans les étoiles et sur Terre obéit aux mêmes lois physiques. Les proportions d’éléments varient cependant selon les milieux. La présence d’un élément dans la croûte terrestre ne garantit pas un gisement facile à exploiter. Extraire, recycler et utiliser des métaux soulève des questions de ressources, d’énergie et d’environnement. Pour vérifier un document, identifier l’auteur, la date, l’unité et le sens du pourcentage.\par
\begin{pcEncadre}{method}{Un exemple expliqué}À l’échelle $1\unite{cm}$ pour $10$ millions de kilomètres, une distance de $150$ millions de kilomètres devient $15\unite{cm}$. Une distance de $750$ millions devient $75\unite{cm}$. Cette échelle des distances ne permet pas de dessiner simultanément des planètes lisibles à leur vraie taille relative : il faut annoncer une autre échelle pour leurs diamètres.\end{pcEncadre}
\pcPage{Expérimenter et retenir}{La matière dans l’Univers}
\phantomsection\label{activite-7}
\begin{pcEncadre}{method}{Investigation en classe}\textbf{Mission.} Construire une frise des distances moyennes au Soleil à partir d’un tableau fourni par le professeur. Choisir une échelle compatible avec la salle ; calculer les positions avant de tracer. Représenter les planètes par des symboles de taille fixe, avec la mention « diamètres non à l’échelle ». Comparer un document de vulgarisation à une source d’agence spatiale : distinguer données mesurées, valeurs arrondies et illustration artistique. Préparer une explication orale de la différence entre galaxie et système planétaire.\end{pcEncadre}
\pcSousTitre{Ce que je dois savoir expliquer}\begin{questions}
\item Le Soleil est une étoile, la Terre une planète et la Lune un satellite naturel.
\item Le Système solaire appartient à la Voie lactée.
\item Une année-lumière mesure une distance.
\item Les éléments terrestres s’inscrivent dans une histoire cosmique.
\item Un schéma astronomique doit préciser ses échelles.
\end{questions}
\begin{pcEncadre}{attention}{Une erreur à éviter}Une image montrant les planètes alignées ne signifie pas qu’elles restent alignées : c’est une convention de présentation de leur ordre.\end{pcEncadre}
\pcRepere{cor-activite-7}{Pistes pour l’investigation}\pcGrille
\clearpage\markright{Je démarre}
\rubrique{Je démarre}
\par\noindent\begin{minipage}{\linewidth}
\begin{exercice}[1]{Classer des objets}
\begin{questions}
\item Classer Soleil, Terre et Lune en étoile, planète ou satellite naturel.
\item Le Soleil appartient-il à la Voie lactée ?
\end{questions}
\end{exercice}
\end{minipage}
\par\noindent\begin{minipage}{\linewidth}
\begin{exercice}[1]{Le bon emboîtement}
\begin{questions}
\item Ranger : Système solaire, Terre, Univers observable, Voie lactée.
\item Une galaxie est-elle une planète géante ?
\end{questions}
\end{exercice}
\end{minipage}
\par\noindent\begin{minipage}{\linewidth}
\begin{exercice}[1]{Une unité trompeuse}
\begin{questions}
\item L’année-lumière est-elle une durée ou une distance ?
\item Pourquoi son nom contient-il « année » ?
\end{questions}
\end{exercice}
\end{minipage}
\par\noindent\begin{minipage}{\linewidth}
\begin{exercice}[1]{L’ordre des planètes}
\begin{questions}
\item Quelles planètes encadrent la Terre dans l’ordre des distances moyennes au Soleil ?
\item Nommer la plus éloignée des huit planètes.
\end{questions}
\end{exercice}
\end{minipage}
\par\noindent\begin{minipage}{\linewidth}
\begin{exercice}[1]{Lire une échelle}
\begin{questions}
\item À $1\unite{cm}$ pour $10$ millions de kilomètres, représenter $50$ millions de kilomètres.
\item Représenter $150$ millions de kilomètres.
\end{questions}
\end{exercice}
\end{minipage}
\clearpage\markright{Je m’entraîne}
\rubrique{Je m’entraîne}
\par\noindent\begin{minipage}{\linewidth}
\begin{exercice}[2]{Comparer des distances}
\begin{questions}
\item Deux distances valent $150$ et $600$ millions de kilomètres. Combien de fois la seconde est-elle plus grande ?
\item Quel écart les sépare ?
\end{questions}
\end{exercice}
\end{minipage}
\par\noindent\begin{minipage}{\linewidth}
\begin{exercice}[2]{Des milliards d’années}
\begin{questions}
\item Écrire $4{,}6$ milliards d’années en millions d’années.
\item Est-ce l’âge d’une personne ou une échelle géologique ?
\end{questions}
\end{exercice}
\end{minipage}
\par\noindent\begin{minipage}{\linewidth}
\begin{exercice}[2]{Les mêmes éléments}
\begin{questions}
\item On détecte de l’hydrogène dans une étoile. Cet hydrogène correspond-il à un élément propre à cette étoile ?
\item Pourquoi les proportions peuvent-elles différer de celles de la Terre ?
\end{questions}
\end{exercice}
\end{minipage}
\par\noindent\begin{minipage}{\linewidth}
\begin{exercice}[2]{Les limites d’une affiche}
\begin{questions}
\item Une affiche donne la même taille à toutes les planètes. Est-elle forcément inutile ?
\item Quelle légende doit-elle porter ?
\end{questions}
\end{exercice}
\end{minipage}
\par\noindent\begin{minipage}{\linewidth}
\begin{exercice}[2]{Passer à une autre échelle}
\begin{questions}
\item Une frise utilise $1\unite{cm}$ pour $25$ millions de kilomètres. Placer une planète à $750$ millions de kilomètres.
\item Quelle distance réelle représente $8\unite{cm}$ ?
\end{questions}
\end{exercice}
\end{minipage}
\clearpage\markright{Je raisonne}
\rubrique{Je raisonne}
\par\noindent\begin{minipage}{\linewidth}
\begin{exercice}[3]{Une preuve d’autorité ?}
\begin{questions}
\item Deux sites donnent des valeurs légèrement différentes pour une distance planétaire. Nommer deux raisons possibles autres qu’une erreur.
\item Quelles informations chercher ?
\end{questions}
\end{exercice}
\end{minipage}
\par\noindent\begin{minipage}{\linewidth}
\begin{exercice}[3]{Un métal « rare »}
\begin{questions}
\item Une fiche annonce qu’un élément est présent dans beaucoup de roches. Cela garantit-il une extraction facile ?
\item Que faut-il encore examiner ?
\end{questions}
\end{exercice}
\end{minipage}
\par\noindent\begin{minipage}{\linewidth}
\begin{exercice}[3]{Une autre unité}
\begin{questions}
\item On prend $1\unite{ua}=150$ millions de kilomètres. Convertir $450$ millions de kilomètres en ua.
\item Convertir $5\unite{ua}$ en millions de kilomètres.
\end{questions}
\end{exercice}
\end{minipage}
\par\noindent\begin{minipage}{\linewidth}
\begin{exercice}[3]{Un alignement dessiné}
\begin{questions}
\item Un schéma aligne les huit planètes. Peut-on l’utiliser pour prédire leur position demain ?
\item Que représente-t-il utilement ?
\end{questions}
\end{exercice}
\end{minipage}
\par\noindent\begin{minipage}{\linewidth}
\begin{exercice}[3]{Rédiger une histoire}
\begin{questions}
\item Pourquoi trouve-t-on du carbone à la fois sur Terre et dans certaines étoiles ?
\item Pourquoi ne faut-il pas dire que tout le carbone a été créé sur Terre ?
\end{questions}
\end{exercice}
\end{minipage}
\pcPage{Défi du chapitre}{La matière dans l’Univers}
\phantomsection\label{defi-ex-7}
\pcDefiTitre{7}{La frise impossible ?}
\begin{questions}
\item La distance moyenne maximale retenue est $4\,500$ millions de kilomètres. La salle offre $9\unite m$. Choisir une échelle simple utilisant toute la longueur.
\item Placer la Terre à $150$ millions de kilomètres.
\item Pourquoi ne peut-on pas exiger aussi des diamètres lisibles à la même échelle ?
\end{questions}
\pcGrille
